\problemname{Tea Brewing}
Egon is going to brew lots of tea for $N$ programming olympiad participants.
He has $K$ tea bags, all of different kinds.
Bag $i$ has tea for $x_i$ people.
It is guaranteed that the bags together are enough for at least $N$ people.

Egon plans to use brewing pots that hold tea for a maximum of 10 people.
Since the bags are of different kinds,
it is not possible to mix multiple bags in the same pot.
However, the same bag can be used for multiple pots.
How many pots must Egon use?

\section*{Input}
The first line contains two integers $K$ and $N$ ($1 \le K \le 10$, $1 \le N \le 100$), the number of tea bags Egon has and the number of programming olympiad participants.

The second line contains the $K$ integers $x_i$ ($1 \le x_i \le 100$), the number of people each bag is enough for.

It is guaranteed that the tea bags are always enough for $N$ people.

\section*{Output}
The program should print an integer: the minimum number of tea pots Egon must use.

\section*{Scoring}
Your solution will be tested on a set of test groups.
To earn points for a group, you must pass all test cases in that group.

\noindent
\begin{tabular}{| l | l | l |}
  \hline
  \textbf{Group} & \textbf{Points} & \textbf{Constraints} \\ \hline
  $1$    & $20$        &  $K = 1$\\ \hline
  $2$    & $80$        &  No additional constraints. \\ \hline
\end{tabular}


\section*{Explanation of Samples}
In the first sample, Egon chooses to brew two pots with the first tea bag and two pots with the third tea bag.
This gives $20+17$ cups of tea, which
is enough for the $36$ participants.

In the second sample, the optimal solution is to brew six pots with the first tea bag,
three pots with the third tea bag, and two with the fourth tea bag.
This gives $54+30+16$ cups of tea, which is enough for the $100$ participants.
